\documentclass[11pt,a4paper]{article}
\usepackage[T1]{fontenc}
\usepackage{lmodern}
\usepackage[margin=27mm]{geometry}
\usepackage{amsmath,amssymb,amsthm,mathtools}
\usepackage{microtype}
\usepackage{booktabs,array}
\usepackage{enumitem}
\usepackage[hidelinks]{hyperref}
\setlength{\parskip}{5pt plus 1pt}
\setlength{\parindent}{0pt}
\setlength{\emergencystretch}{3em}
\setlist{itemsep=2pt,topsep=3pt}
\raggedbottom
\newtheorem{theorem}{Theorem}[section]
\newtheorem{lemma}[theorem]{Lemma}
\newtheorem{corollary}[theorem]{Corollary}
\newtheorem{proposition}[theorem]{Proposition}
\theoremstyle{definition}
\newtheorem{definition}[theorem]{Definition}
\newtheorem{example}[theorem]{Example}
\theoremstyle{remark}
\newtheorem{remark}[theorem]{Remark}
\newcommand{\Z}{\mathbb{Z}}
\newcommand{\N}{\mathbb{N}}
\newcommand{\Q}{\mathbb{Q}}
\newcommand{\F}{\mathbb{F}}
\newcommand{\ps}[1]{#1[[x]]}
\newcommand{\coeff}[2]{[x^{#1}]\,#2}
\newcommand{\hl}[1]{\textbf{#1}}
\newcommand{\para}[1]{\par\medskip\noindent\textbf{#1}\enspace}
\hypersetup{pdftitle={Congruences for the series A with A(x - x A(7x)/7) = x},pdfauthor={Dev D. Goyal},pdfsubject={OEIS A393859}}
\title{Congruences for the series $A$ with $A\bigl(x-x\,A(7x)/7\bigr)=x$\\[6pt]\large A proof of a conjecture in OEIS A393859}
\author{Dev D. Goyal}
\date{October 2026}
\begin{document}
\maketitle
\begin{abstract}
The sequence A393859 of the On-Line Encyclopedia of Integer Sequences (OEIS) lists the coefficients of the power
series $A(x)=x+x^2+9x^3+481x^4+\cdots$ defined by the equation
$A\bigl(x-x\,A(7x)/7\bigr)=x$. The unknown series $A$ appears on both sides of this equation, and its coefficients grow very fast.
Paul~D.~Hanna conjectured that every coefficient leaves remainder $1$ on division by $8$. We prove this. In fact we prove that, with $7$ replaced by any positive integer $q$, the series becomes $x+x^2+x^3+\cdots$ once its coefficients are reduced modulo $q+1$. The proof rests on the congruence $q\equiv-1\pmod{q+1}$: for $q=-1$ the equation
has an explicit solution. The results have also been checked by computer and proved formally in the
Lean proof assistant.
\end{abstract}
\begin{center}\begin{minipage}{0.88\linewidth}\small
\textbf{How this paper was produced.} This paper reports the result of an AI research effort directed by
Dev~D.~Goyal. ChatGPT (OpenAI), using the GPT-6 model with Pro reasoning and working autonomously from prompts by
Dev~D.~Goyal, found the proof, wrote its Lean~4 formalization, and wrote a first draft of the paper.
Dev~D.~Goyal chose the problems and asked for proofs, formal verification and revisions, but did not supply proof
ideas. Claude Opus~5.5, used with Claude Code, then checked the work independently: it rebuilt the formal proofs from
source and confirmed that they rely only on Lean's three standard axioms, recomputed the sequence from its definition
and tested every claim against the published data, and read every step of the proofs. It also wrote the present
version of this paper.
Dev~D.~Goyal has verified the results and is responsible for them. The paper has not been peer reviewed.
The section ``About this paper'' at the end gives the details.
\end{minipage}\end{center}

\section{Introduction}

\subsection{The sequence}\label{sec:sequence}
The \emph{On-Line Encyclopedia of Integer Sequences} (OEIS)~\cite{oeisdb} is a large public database of integer
sequences. Many
entries also record \emph{conjectures}: patterns that the contributor observed in the data but could not prove.

This paper is about entry A393859~\cite{oeis}, contributed by Paul~D.~Hanna in March 2026. It defines a power series
\[
A(x)=a(1)\,x+a(2)\,x^2+a(3)\,x^3+\cdots
\]
by the equation
\begin{equation}\label{eq:oeis}
A\bigl(x-x\,A(7x)/7\bigr)=x.
\end{equation}
Here $A(7x)$ is the series $A$ with $x$ replaced by $7x$, so its coefficient of $x^n$ is
$7^n\,a(n)$. Dividing by $7$ gives the series whose coefficients are
$7^{\,n-1}a(n)$, which are still integers. The left side of \eqref{eq:oeis} means: substitute the series
$x-x\,A(7x)/7$ into $A$. (Section~\ref{sec:background} explains what substituting one
series into another means.) So the equation says that $A$ undoes a substitution that is itself built from $A$. Because
$A$ appears on both sides, it is not obvious that the equation has a solution, or only one; Theorem~\ref{thm:exist}
shows that it has exactly one solution of the form $x+(\text{higher powers of }x)$.

The first terms, starting at $n=1$, are 1, 1, 9, 481, 167937, 404407577, 6800120127457, \dots{} They grow very fast: $a(15)$
already has 78 digits. The entry reads:
\begin{quote}\small\ttfamily\raggedright G.f. A(x) satisfies A( x - x*A(7*x)/7 ) = x.\par\smallskip Conjecture: a(n) == 1 (mod 8) for n >= 1.\par\smallskip\hfill\textrm{--- OEIS A393859, Paul D. Hanna, March 2026}\end{quote}

Here ``G.f.''\ stands for \emph{generating function}, the series whose coefficients are the terms of the sequence.
The entry also records other formulas for $A(x)$, for example as an infinitely nested sequence of substitutions; we do
not need them.

\subsection{The question}
Hanna conjectured, from the computed terms, that every $a(n)$ leaves remainder $1$ on division by
$8$. The table shows the remainders of the first $32$ terms modulo $8$.
\begin{center}\small\setlength{\tabcolsep}{3.2pt}
\begin{tabular}{>{\raggedright\arraybackslash}p{6.4em}>{\centering\arraybackslash}p{1.15em}>{\centering\arraybackslash}p{1.15em}>{\centering\arraybackslash}p{1.15em}>{\centering\arraybackslash}p{1.15em}>{\centering\arraybackslash}p{1.15em}>{\centering\arraybackslash}p{1.15em}>{\centering\arraybackslash}p{1.15em}>{\centering\arraybackslash}p{1.15em}>{\centering\arraybackslash}p{1.15em}>{\centering\arraybackslash}p{1.15em}>{\centering\arraybackslash}p{1.15em}>{\centering\arraybackslash}p{1.15em}>{\centering\arraybackslash}p{1.15em}>{\centering\arraybackslash}p{1.15em}>{\centering\arraybackslash}p{1.15em}>{\centering\arraybackslash}p{1.15em}}
\toprule
$n$ & 1 & 2 & 3 & 4 & 5 & 6 & 7 & 8 & 9 & 10 & 11 & 12 & 13 & 14 & 15 & 16 \\ \midrule
$a(n)\bmod 8$ & 1 & 1 & 1 & 1 & 1 & 1 & 1 & 1 & 1 & 1 & 1 & 1 & 1 & 1 & 1 & 1 \\
\bottomrule
\end{tabular}
\\[4pt]
\begin{tabular}{>{\raggedright\arraybackslash}p{6.4em}>{\centering\arraybackslash}p{1.15em}>{\centering\arraybackslash}p{1.15em}>{\centering\arraybackslash}p{1.15em}>{\centering\arraybackslash}p{1.15em}>{\centering\arraybackslash}p{1.15em}>{\centering\arraybackslash}p{1.15em}>{\centering\arraybackslash}p{1.15em}>{\centering\arraybackslash}p{1.15em}>{\centering\arraybackslash}p{1.15em}>{\centering\arraybackslash}p{1.15em}>{\centering\arraybackslash}p{1.15em}>{\centering\arraybackslash}p{1.15em}>{\centering\arraybackslash}p{1.15em}>{\centering\arraybackslash}p{1.15em}>{\centering\arraybackslash}p{1.15em}>{\centering\arraybackslash}p{1.15em}}
\toprule
$n$ & 17 & 18 & 19 & 20 & 21 & 22 & 23 & 24 & 25 & 26 & 27 & 28 & 29 & 30 & 31 & 32 \\ \midrule
$a(n)\bmod 8$ & 1 & 1 & 1 & 1 & 1 & 1 & 1 & 1 & 1 & 1 & 1 & 1 & 1 & 1 & 1 & 1 \\
\bottomrule
\end{tabular}
\end{center}

\subsection{What this paper proves}
\textbf{Main result} (Corollary~\ref{cor:oeis}). \emph{The conjecture is true: $a(n)\equiv1\pmod{8}$ for every $n\ge1$.}

The number $7$ plays no special role. Hanna contributed a family of such sequences, in which
$7$ is replaced by another positive integer $q$, and we treat all of them at once.
Theorem~\ref{thm:mod} shows that, for every $q\ge1$, all coefficients leave remainder $1$ on division by $q+1$.
Theorem~\ref{thm:parity} shows that, for every even $q$, a coefficient is odd exactly when its index is a power of $2$.

The proof rests on two observations. Modulo $q+1$ we have $q\equiv-1$, and for $q=-1$ the equation has the
explicit solution $x+x^2+x^3+\cdots$. Similarly, modulo $2$ an even $q$ is congruent to $0$, and for $q=0$ the
equation becomes the classical equation of the Catalan numbers, whose parity is well known.

\subsection{What was known}
Hanna contributed the four entries A393856--A393859, for $q=4,5,6,7$, in March 2026, each with the conjectures for its
own value of $q$. The OEIS also lists the members for $q=0,1,2,3$, which were contributed earlier. For $q=0$, where
$A(qx)/q$ is replaced by $x$ (its value at $q=0$ according to formula \eqref{eq:L} below), the equation becomes
$A(x-x^2)=x$; its solution is the generating series of the Catalan numbers (A000108). The cases $q=1,2,3$ are
A088714, A195195 and A195196. We are not aware of an earlier proof of the conjectures. The tools we use, formal power
series and their reduction modulo an integer, are standard; see~\cite{niven,stanley}.

\subsection{Why it is interesting}
Equations in which an unknown series is substituted into an expression built from itself are common in
combinatorics, where they describe structures defined recursively, such as trees. The Catalan numbers, the case $q=0$
of this family, are the classic example. Hanna has contributed many such equations to the OEIS. For $q\ge1$ the
coefficients grow rapidly and no closed formula for them is known to us, yet their remainders follow simple rules. The
proof explains these rules: working modulo $q+1$ replaces the parameter $q$ by $-1$, and, for even $q$, working modulo
$2$ replaces it by $0$; at these two values the equation can be solved explicitly. The same approach, reduction modulo
a number for which the equation simplifies, is used in other papers of this collection.

\subsection{How to read this paper}
Section~\ref{sec:background} explains the few notions we need: power series, substitution, and remainders.
Section~\ref{sec:results} makes the equation precise for every $q$ and states the results, and
Section~\ref{sec:proofs} proves them. Section~\ref{sec:checks} describes how the results were checked, by computer
and by a formal proof in Lean. Section~\ref{sec:discussion} discusses the scope of the results, a further congruence,
and open questions.
\section{Background}\label{sec:background}
This section collects the definitions and facts used later. Readers who know formal power series and congruences can skip it and return when needed.
\subsection{Formal power series}
A \emph{formal power series} is an expression
\[
U(x)=u_0+u_1x+u_2x^2+u_3x^3+\cdots ,
\]
determined by its sequence of coefficients $u_0,u_1,u_2,\dots$; see~\cite{niven,stanley}. No number is ever
substituted for $x$, so questions of convergence do not arise. We write $\coeff{n}{U}=u_n$ for the coefficient of $x^n$ in $U$; two series are equal
when all their coefficients are equal.

Series are added coefficient by coefficient and multiplied like polynomials:
\[
\coeff{n}{(UV)}=u_0v_n+u_1v_{n-1}+\cdots+u_nv_0 .
\]
Each coefficient of a product is a finite sum. For example, $(1-x)(1+x+x^2+x^3+\cdots)=1$, so $1+x+x^2+\cdots$ is
the \emph{multiplicative inverse} of $1-x$, written $\frac{1}{1-x}$. More generally, every series $D$ with constant term
$1$ has a unique multiplicative inverse $\frac1D$, whose coefficients are determined one at a time by $D\cdot\frac1D=1$
using only addition, subtraction and multiplication. Whenever a quotient of power series appears in this paper, its
denominator has constant term $1$, so the quotient is again a power series.

For a commutative ring $R$, such as $\Z$, $\Q$ or the integers modulo $m$, we write $\ps{R}$ for the set of power
series with coefficients in $R$; for instance $\ps{\F_2}$ consists of the series with coefficients modulo $2$. We write
$U(0)$ for the constant term $u_0$ of $U$, and $U=O(x^d)$ when the coefficients of $1,x,\dots,x^{d-1}$ in $U$ are all
zero. Thus $U-V=O(x^d)$ means that $U$ and $V$ have the same coefficient of $x^n$ for every $n<d$. A series with
constant term $0$ and coefficient of $x$ equal to $1$, that is, a series of the form $x+u_2x^2+u_3x^3+\cdots$, is
called \emph{normalized}.

\subsection{Substituting one series into another}
If $V(x)$ has constant term $0$, we can substitute $V$ for $x$ in $U$:
\[
U(V(x))=u_0+u_1V(x)+u_2V(x)^2+u_3V(x)^3+\cdots .
\]
This is called \emph{composition}, also written $U\circ V$. It makes sense because $V(x)^k$ has no terms below $x^k$,
so only the finitely many terms with $k\le n$ affect the coefficient of $x^n$. For example, if $U=x+x^2$ and $V=x^2$, then
$U(V(x))=x^2+x^4$. Composition is associative: $U\circ(V\circ W)=(U\circ V)\circ W$.

Substituting a series into itself repeatedly gives its \emph{iterates}: $U^{\circ2}(x)=U(U(x))$,
$U^{\circ3}(x)=U(U(U(x)))$, and so on, with $U^{\circ0}(x)=x$ and $U^{\circ1}=U$. They must not be confused with
\emph{powers} $U^2=U\cdot U$, $U^3=U\cdot U\cdot U$. (The OEIS sometimes writes $A^k(x)$ for the $k$-th iterate. In
this paper the small circle always means an iterate, and a plain exponent always means a power.)

\subsection{Undoing a substitution}
A normalized series can always be ``undone'' by another normalized series.

\begin{lemma}[Compositional inverse]\label{lem:inverse}
For every normalized series $U$ there is exactly one normalized series $V$ with $U(V(x))=x$, and this $V$ also
satisfies $V(U(x))=x$. If $U$ has integer coefficients, so does $V$.
\end{lemma}

We call $V$ the \emph{compositional inverse} of $U$. The OEIS calls it the \emph{series reversion} of $U$. For example,
the compositional inverse of $x-x^2$ is $C(x)=x+x^2+2x^3+5x^4+14x^5+\cdots$, whose coefficients are the Catalan numbers.
Indeed $C-C^2=x$, which says exactly that substituting $C$ into $x-x^2$ gives $x$.

\begin{proof}
Write $V=x+v_2x^2+v_3x^3+\cdots$ with unknown coefficients. In
$U(V)=V+u_2V^2+u_3V^3+\cdots$, the coefficient of $x^n$ (for $n\ge2$) equals $v_n$ plus terms that involve only
$v_2,\dots,v_{n-1}$. Setting it equal to $0$ therefore determines $v_n$ from the earlier coefficients, using only
additions, subtractions and multiplications. So there is exactly one $V$ with $U\circ V=x$, and it has integer
coefficients if $U$ does. Applying the same construction to $V$ gives a normalized $W$ with $V\circ W=x$. Then
$U=U\circ(V\circ W)=(U\circ V)\circ W=W$, so $V\circ U=V\circ W=x$ as well.
\end{proof}

\subsection{Remainders and congruences}
For an integer $m\ge2$ we write $a\equiv b\pmod m$ (``$a$ is congruent to $b$ modulo $m$'') when $a-b$ is divisible by
$m$, that is, when $a$ and $b$ leave the same remainder after division by $m$. The remainder of $a$ on division by
$m$ is called its \emph{residue} modulo $m$; it is one of $0,1,\dots,m-1$.

Addition and multiplication are compatible with congruences, so the residues $0,1,\dots,m-1$, added and multiplied
modulo $m$, form a ring, written $\Z/m\Z$. When $m=p$ is a prime we also write $\F_p$. In $\F_2=\{0,1\}$ we have $1+1=0$,
so addition and subtraction coincide.

For power series with integer coefficients, $U\equiv V\pmod m$ means that \emph{every} coefficient of $U-V$ is
divisible by $m$. Equivalently, the two series become equal when every coefficient is replaced by its residue; we call
this \emph{reducing} the series modulo $m$. The key fact is the following.

\emph{Reduction modulo $m$ is compatible with all the operations above.} To reduce a sum, a product, a substitution, or
an inverse, we may reduce the ingredients first and then compute in $\Z/m\Z$. This holds because every
coefficient of the result is obtained from finitely many coefficients of the ingredients by adding, subtracting and
multiplying, and these operations are compatible with remainders.

\para{Rings.} Some statements below hold over every \emph{commutative ring}, that is, every set with an addition,
subtraction and multiplication satisfying the usual rules of arithmetic (for example $\Z$, $\Q$ and $\Z/m\Z$). No
division is used; this matters because in $\Z/m\Z$ division is not always possible (for instance, $2$ has no inverse in
$\Z/4\Z$).

\subsection{The method}
Almost every proof in this paper follows the same three steps.
\begin{enumerate}
\item \emph{Uniqueness.} Show that the defining equation determines the coefficients one at a time, using only
addition, subtraction and multiplication. Then the equation has exactly one solution of the required shape, in every
commutative ring, including $\Z/m\Z$.
\item \emph{Reduction.} Reducing the integer solution modulo $m$ gives a solution of the reduced equation,
because reduction is compatible with all operations.
\item \emph{A candidate solution.} Find an explicit series that satisfies the reduced equation. By uniqueness, it is
the reduction of the sequence, and its coefficients can be read off.
\end{enumerate}
Finding the candidate in step~3 is usually the main difficulty.

\section{Results}\label{sec:results}

To treat the whole family at once, we replace $7$ by a parameter $q$. First we make the equation
precise. For a series $U=u_1x+u_2x^2+u_3x^3+\cdots$ without constant term and a number $q\ne0$, the series $U(qx)/q$
has the coefficients $q^{\,n-1}u_n$. We take this formula as the definition, for every $q$:
\begin{equation}\label{eq:L}
L_q(U)=\sum_{n\ge1}q^{\,n-1}u_n\,x^n=u_1x+q\,u_2x^2+q^2u_3x^3+\cdots .
\end{equation}
No division by $q$ occurs, so $L_q(U)$ makes sense for every $q$, including $q=0$ (where it equals $u_1x$), and for
coefficients in every commutative ring, such as $\Z/m\Z$. The equation of the family is
\begin{equation}\label{eq:def}
A\bigl(x-x\,L_q(A)\bigr)=x,\qquad A\ \text{normalized}.
\end{equation}
For $q=7$ this is equation \eqref{eq:oeis} of the entry.

\begin{theorem}[Exactly one solution]\label{thm:exist}
Let $q$ be an element of a commutative ring $R$. There is exactly one normalized series $G$ with coefficients in $R$
that satisfies \eqref{eq:def}. When $R=\Z$ and $q$ is an integer, we call this solution $G_q$. For every $m\ge2$, the
reduction of $G_q$ modulo $m$ is the solution of \eqref{eq:def} over $\Z/m\Z$, with $q$ replaced by its residue.
\end{theorem}

So the sequence of OEIS A393859 is $a(n)=\coeff{n}{G_{7}}$.

\begin{theorem}[All coefficients are $1$ modulo $q+1$]\label{thm:mod}
For every integer $q\ge1$,
\[
G_q(x)\equiv\frac{x}{1-x}=x+x^2+x^3+\cdots\pmod{q+1}.
\]
That is, every coefficient $\coeff{n}{G_q}$ with $n\ge1$ leaves remainder $1$ on division by $q+1$.
\end{theorem}

\begin{theorem}[Parity]\label{thm:parity}
If $q$ is even, then
\[
G_q(x)\equiv x+x^2+x^4+x^8+\cdots=\sum_{k\ge0}x^{2^k}\pmod 2,
\]
so $\coeff{n}{G_q}$ is odd if and only if $n$ is a power of $2$. If $q$ is odd, every $\coeff{n}{G_q}$ with $n\ge1$
is odd.
\end{theorem}

\begin{corollary}[Hanna's conjecture]\label{cor:oeis}
For the sequence of OEIS A393859: $a(n)\equiv1\pmod{8}$ for every $n\ge1$.
\end{corollary}

\section{Proofs}\label{sec:proofs}

\para{The idea in brief.} The equation determines the coefficients of $A$ one at a time, without any division, so it
has exactly one solution over every commutative ring, including $\Z/(q+1)\Z$ and $\Z/2\Z$
(Theorem~\ref{thm:exist}). Modulo $q+1$ we may replace $q$ by $-1$, and then a direct calculation shows that $x+x^2+x^3+\cdots$ is
a solution. Modulo $2$, for even $q$, we may replace $q$ by $0$, and then the equation says that $A$ undoes the
substitution $x\mapsto x-x^2$, which is exactly what the Catalan series does. By uniqueness, these explicit series
are the reductions of $G_q$.

\begin{proof}[Proof of Theorem~\ref{thm:exist}]
Write $G=x+g_2x^2+g_3x^3+\cdots$ with unknown coefficients, and put $V=x-x\,L_q(G)$. By \eqref{eq:L}, the coefficient
$g_k$ enters $L_q(G)$ only in degree $k$, and therefore enters $V$ only in degree $k+1$, because of the factor $x$. Also,
$V=x+O(x^2)$ is normalized. Now look at the coefficient of $x^n$, for some $n\ge2$, in
\[
G(V)=V+g_2V^2+g_3V^3+\cdots+g_nV^n+\cdots .
\]
\begin{itemize}
\item The term $g_nV^n$ contributes exactly $g_n$, because $V^n=x^n+O(x^{n+1})$.
\item The terms $g_kV^k$ with $k>n$ contribute nothing, because $V^k$ has no terms below $x^k$.
\item The remaining terms, $V$ and $g_kV^k$ with $2\le k<n$, involve $g_k$ with $k<n$, and the coefficients of $V$ up
to degree $n$, which involve only $g_2,\dots,g_{n-1}$.
\end{itemize}
So the coefficient of $x^n$ in $G(V)$ is $g_n$ plus an expression in $q,g_2,\dots,g_{n-1}$ built with additions,
subtractions and multiplications. Setting it equal to the coefficient of $x^n$ in $x$, which is $0$, determines $g_n$.
Doing this for $n=2,3,4,\dots$ in turn gives exactly one solution, in every commutative ring.

Reduction modulo $m$ is compatible with all the operations in \eqref{eq:def}: with $L_q$, which multiplies the
coefficients by powers of $q$, with products, and with substitution. So the reduction of $G_q$ satisfies the reduced
equation, and by the uniqueness just proved (over $\Z/m\Z$) it is its solution.
\end{proof}

\begin{proof}[Proof of Theorem~\ref{thm:mod}]
We work over $\Z/(q+1)\Z$, where $q=-1$. Let $H=x/(1-x)=x+x^2+x^3+\cdots$. All coefficients of $H$ (except the
constant term) are $1$, so by \eqref{eq:L}
\[
L_{-1}(H)=x-x^2+x^3-x^4+\cdots=\frac{x}{1+x},
\qquad\text{hence}\qquad
x-x\,L_{-1}(H)=x-\frac{x^2}{1+x}=\frac{x}{1+x}.
\]
Substituting this into $H$ gives
\[
H\!\left(\frac{x}{1+x}\right)=\frac{x/(1+x)}{1-x/(1+x)}=\frac{x/(1+x)}{1/(1+x)}=x .
\]
So $H$ satisfies the reduced equation. By Theorem~\ref{thm:exist}, $H$ is the reduction of $G_q$ modulo $q+1$. Every
coefficient of $H$ of positive degree is $1$.
\end{proof}

\begin{proof}[Proof of Theorem~\ref{thm:parity}]
If $q$ is odd, then $q+1$ is even. By Theorem~\ref{thm:mod}, every coefficient of positive degree is $1$ plus a multiple
of $q+1$, hence odd.

Now let $q$ be even, and work over $\F_2$, where $q=0$. By \eqref{eq:L}, $L_0(U)=u_1x$, which is $x$ for a normalized
$U$, so over $\F_2$ the equation becomes $G(x-x^2)=x$. Thus $G$ is the compositional inverse of $x-x^2$
(Lemma~\ref{lem:inverse}), which is the reduction of the Catalan series $C$ (Section~\ref{sec:background}). It is a
classical fact that $C\equiv x+x^2+x^4+x^8+\cdots\pmod 2$: indeed, $S=\sum_kx^{2^k}$ satisfies $S-S^2=x$ over $\F_2$,
because squaring a series over $\F_2$ doubles all its exponents (there $(a+b)^2=a^2+2ab+b^2=a^2+b^2$ and $c^2=c$), so $S^2=\sum_{k\ge1}x^{2^k}$; and the equation
$U-U^2=x$ determines $U$ uniquely, one coefficient at a time. Hence $G_q\equiv S\pmod2$.
\end{proof}

Corollary~\ref{cor:oeis} is the case $q=7$ of Theorems~\ref{thm:mod} and~\ref{thm:parity}.
(For odd $q$ the parity statement adds nothing new: it already follows from the congruence modulo $q+1$. Since $8=2\cdot4$, the corollary also says that every $a(n)$ is odd and leaves remainder $1$ on division by $4$.)

\section{Checking the results}\label{sec:checks}

A mathematical proof can contain mistakes, so the results were checked in two independent ways.

\para{Numerical check.}
We solved \eqref{eq:oeis} coefficient by coefficient with exact integer arithmetic, up to $n=25$; this
reproduces all 13 terms printed in the OEIS entry. We then ran the same computation
modulo $16=2(q+1)$, which is legitimate because the construction uses only additions,
subtractions and multiplications (Theorem~\ref{thm:exist}). It agrees with the exact values for $n\le25$, and it confirms
the congruence of Corollary~\ref{cor:oeis} for every $n<300$.
The computation was done with the PARI/GP computer algebra system, using code that was written separately from the
formal proofs (file \texttt{verification/independent\_checks.gp} in the accompanying repository~\cite{repo}).
A computation can only test finitely many cases, so it is a check on the definitions and on the indexing, not a proof.

\para{Formal proof in Lean.}
Lean~\cite{lean} is a \emph{proof assistant}: a program in which mathematical statements and their proofs are written
in a precise formal language, and every step is checked mechanically down to the basic rules of logic. The theorems of Section~\ref{sec:results} are proved in Lean~4 (version 4.19.0) on top of
Mathlib~\cite{mathlib}, Lean's library of formalized mathematics; any step proved only on paper is stated below. The statements are collected in the file \texttt{lean/Entries/A393859.lean} of the repository,
in the namespace \texttt{OEIS.A393859}. In Lean, the operation $L_q$ of \eqref{eq:L} is called \texttt{scaleTail q}. The general theorems for every $q$ are \texttt{integer\_mod\_parameter\_add\_one} (Theorem~\ref{thm:mod}) and \texttt{integer\_coeff\_odd\_iff} (Theorem~\ref{thm:parity} for even $q$), in the namespace \texttt{Hanna.\allowbreak Related.\allowbreak ScaledInverse}, in the files \texttt{lean/\allowbreak Hanna/\allowbreak Related/\allowbreak ScaledInverse.lean} and \texttt{ScaledParity.lean} of the same folder.
The table lists where each result of this paper is proved.
\begin{center}\small\renewcommand{\arraystretch}{1.25}
\begin{tabular}{>{\raggedright\arraybackslash}p{0.52\linewidth}>{\raggedright\arraybackslash}p{0.4\linewidth}}
\toprule
Statement in this paper & Lean theorem(s)\\
\midrule
Theorem~\ref{thm:exist}, over every commutative ring & \texttt{source\_\allowbreak{}existence\_\allowbreak{}uniqueness} \\
$G_{7}$ is normalized and satisfies \eqref{eq:oeis} & \texttt{source\_\allowbreak{}normalized}, \texttt{source\_\allowbreak{}equation} \\
Corollary~\ref{cor:oeis}: $a(n)\equiv1\pmod{8}$ for $n\ge1$ & \texttt{a393859\_\allowbreak{}mod\_\allowbreak{}eight} \\
\bottomrule
\end{tabular}
\end{center}
The Lean project was rebuilt from source, and the axioms on which each of these theorems depends were listed with
Lean's \texttt{\#print axioms} command. Every one depends only on the three standard axioms of Lean's logic
(\texttt{propext}, \texttt{Classical.choice}, \texttt{Quot.sound}); no step is left unproved. A formal proof certifies
exactly the statement that Lean checks. Whether that statement says what the OEIS entry means is a separate question,
so the Lean definitions were also compared with the definition in the OEIS entry
(Section~\ref{sec:sequence}).

\section{Discussion}\label{sec:discussion}

\para{Scope.} Corollary~\ref{cor:oeis} settles the conjecture of the entry completely,
for every index $n$. The theorems describe the remainders modulo $q+1$ and modulo $2$; they say nothing about the size
of the coefficients.

\para{A further congruence.} The argument used for Theorem~\ref{thm:parity} works modulo $q$ itself. In arithmetic
modulo $q$ the parameter $q$ is $0$, so the equation becomes $G(x-x^2)=x$ again, and $G_q$ agrees modulo $q$ with the
Catalan series:
\[
\coeff{n}{G_q}\equiv\frac{1}{n}\binom{2n-2}{n-1}\pmod q\qquad(n\ge1).
\]
For example, for $q=7$ and $n=5$, the term $a(5)=167937$ and the Catalan number $\frac15\binom84=14$ both leave remainder $0$ on division by $7$. Since $q$ and $q+1$ have no common factor, a number is determined modulo $q(q+1)$ by its
remainders modulo $q$ and modulo $q+1$ (the Chinese remainder theorem). So, together with Theorem~\ref{thm:mod}, this
determines every coefficient of $G_q$ modulo $q(q+1)=56$ when $q=7$. We checked
this congruence by computer for $q=2,3,\dots,7$ and every $n<300$. It is not part of the Lean formalization.

\para{The rest of the family.} The theorems hold for every $q\ge1$. The OEIS lists the members $q=1,2,3$ as A088714,
A195195 and A195196; we checked that their published terms agree with the coefficients of $G_1$, $G_2$ and $G_3$ (for
A088714 after shifting the index by one). For instance, A195195 is defined by the case $q=2$ of
\eqref{eq:oeis}, so our results show that all its terms leave remainder $1$ on division by $3$, and that a term is odd
exactly when its index is a power of $2$.

\para{Open questions.} We have not studied the remainders modulo other numbers, such as $64=(q+1)^2$
for $q=7$. It would be interesting to know whether they are also given by an explicit series.

\section*{About this paper}\label{sec:about}
\para{How the work was done.}
This paper is one of 29 papers on conjectures from the OEIS that came out of an AI research effort directed by
Dev~D.~Goyal~\cite{repo}. The roles were as follows.
\begin{itemize}
\item \emph{Dev~D.~Goyal} started the project, chose which problems to pursue, asked for proofs, formal verification and
revisions, and decided what to publish. Dev~D.~Goyal did not supply proof strategies or intermediate arguments.
\item \emph{ChatGPT} (OpenAI), using the GPT-6 model with Pro reasoning, worked autonomously from those prompts. It
searched the OEIS for suitable conjectures, found the proofs, wrote the Lean~4 formalization, and wrote a first draft
of this paper.
\item \emph{Claude Opus~5.5}, used with Claude Code, checked the work independently. It rebuilt the Lean project from source and confirmed that the theorems listed in Section~\ref{sec:checks} rely only on
Lean's standard axioms; compared the Lean definitions with the OEIS definition; recomputed the sequence with separately
written PARI/GP code, compared it with the terms published in the OEIS, and tested every claim of the paper on many
terms; and read every step of the proofs. It then wrote the present version of the paper, expanding the
explanations and, in places, writing out an argument in more detail or by a more elementary route. The theorems of
Section~\ref{sec:results} and the proof strategies are ChatGPT's. The Discussion section, including the short arguments
in it, was written by Claude Opus~5.5; its observations were checked by computer and are not part of the Lean
formalization.
\end{itemize}
Dev~D.~Goyal has verified the results and is responsible for this paper.

\para{Review status.} This paper has not been peer reviewed, and it has not been submitted to a journal or to the OEIS.

\para{Priority.} An OEIS entry that labels a statement as a conjecture shows only that its contributor did not know
a proof; it does not show that no proof existed. During the research only limited searches for earlier work were
made. No claim of priority is made.

\para{Corrections.} Please report any error in the issue tracker of the repository~\cite{repo}, where the Lean sources and the numerical checks are also
available.

\section*{Acknowledgements}
The sequence and the conjecture are due to Paul~D.~Hanna~\cite{oeis}, who contributed
the entries A393856--A393859 ($q=4,5,6,7$), as well as the earlier members A088714, A195195 and A195196.

\begin{thebibliography}{99}\small
\raggedright
\bibitem{oeis} P.~D.~Hanna, entry A393859 (G.f.\ $A(x)$ satisfies $A\bigl(x-x\,A(7x)/7\bigr)=x$), \emph{The On-Line Encyclopedia of Integer Sequences}, March 2026. \url{https://oeis.org/A393859}. Text as of October 2026.
\bibitem{oeisdb} N.~J.~A.~Sloane and The OEIS Foundation Inc., \emph{The On-Line Encyclopedia of Integer Sequences}. \url{https://oeis.org}.
\bibitem{stanley} R.~P.~Stanley, \emph{Enumerative Combinatorics}, Volume~2, Cambridge University Press, 1999. (Chapter~5 treats formal power series, composition and compositional inverses.)
\bibitem{niven} I.~Niven, Formal power series, \emph{American Mathematical Monthly} 76 (1969), 871--889.
\bibitem{lean} L.~de~Moura and S.~Ullrich, The Lean~4 theorem prover and programming language, in \emph{Automated Deduction -- CADE 28}, Lecture Notes in Computer Science 12699, Springer, 2021, 625--635.
\bibitem{mathlib} The mathlib Community, The Lean mathematical library, in \emph{Proceedings of the 9th ACM SIGPLAN International Conference on Certified Programs and Proofs (CPP 2020)}, 367--381. Version used: commit \texttt{c44e0c8ee63c}.
\bibitem{repo} D.~D.~Goyal, \emph{Proofs and corrections of OEIS conjectures, with Lean~4 formalizations}, repository, 2026. \url{https://github.com/dvgyl/oeis-a396843-proof}.
\end{thebibliography}
\end{document}
